\section{Complex energies and rational transfer products}\label{app:complex}
This appendix proves the extension of Paper~B, Lemmas~2.4--2.5, used below. It does not assert that the complex-energy Lyapunov exponent equals its real spectral value.

For $z\in\C$, put
\[
 A(z,t)=\begin{pmatrix}z-v(t)&-1\\1&0\end{pmatrix},\qquad
 A_r^{(\omega)}(z,t)=A(z,t+(r-1)\omega)\cdots A(z,t).
\]
On the spectrum of the supercritical almost Mathieu operator, the Lyapunov exponent equals $L=\log\lambda$, by Bourgain--Jitomirskaya~\cite[Corollary~2]{BourgainJitomirskaya2002}, after matching the coupling convention. This spectral identity is the AMO instance of the growth input in the companion papers.

\begin{lemma}[Uniform complex neighbourhood]\label{lem:complexgrowth}
For each $h>L$ there are a bounded open neighbourhood $U_h^{\C}$ of $\Sig$ and $C_h<\infty$ such that
\[
 \norm{A_r^{(\alpha)}(z,t)}\le C_he^{hr}
 \quad(z\in U_h^{\C},\ t\in\T,\ r\ge0).
\]
After enlarging $C_h$, every sufficiently large convergent $p/q$ satisfies
\begin{equation}\label{eq:complexrat}
 \norm{A_r^{(p/q)}(z,t)}\le C_he^{hr}
 \quad(0\le r\le q, z\in U_h^{\C},\ t\in\T).
\end{equation}
The same estimates hold for subblocks.
\end{lemma}
\begin{proof}
Choose $h_1\in(L,h)$. At each $E_0\in\Sig$, the infimum formula for the Lyapunov exponent gives an integer $r_0$ with
\[
 \frac1{r_0}\int_\T\log\norm{A_{r_0}^{(\alpha)}(E_0,t)}\,dt<h_1.
\]
The finite product is jointly continuous in complex $z$ and real $t$. Its determinant is one, so its norm never vanishes. The strict averaged inequality persists on a compact complex neighbourhood $K_0$ of $E_0$.

Rotation by $r_0\alpha$ is uniquely ergodic. The continuous functions $t\mapsto\log\norm{A_{r_0}^{(\alpha)}(z,t)}$, $z\in K_0$, form a compact family in the uniform norm. A finite uniform approximation to that family makes the Birkhoff convergence uniform in both variables. For all sufficiently many blocks, their summed logarithmic norms are at most $mr_0h$. Splitting $r=mr_0+s$ leaves a uniformly bounded remainder. This proves the bound on $K_0$. A finite cover of $\Sig$ gives $U_h^{\C}$ and one constant.

For the rational comparison, each one-step error has norm at most $d=\Lip(v)|\delta|$. A term with $j$ such errors in a product of $r\le q$ factors contains at most $j+1$ irrational subblocks. Its total norm is bounded by $C_h^{j+1}e^{h(r-j)}d^j$. Consequently
\[
 \norm{A_r^{(p/q)}}\le C_he^{hr}(1+C_hde^{-h})^r
 \le C_he^{hr}\exp(C_h\Lip(v)e^{-h}),
\]
because $q|\delta|<1$ for convergents. No step uses realness of $z$. Starting phases cover all subblocks.
\end{proof}

\begin{corollary}[Polynomial derivatives and the top contour]\label{cor:derivatives}
For fixed $r\ge1$, on a fixed smaller complex neighbourhood,
\[
 |\partial_z^r\Delta_t(z)|\le C_{r,h}q^re^{hq},\qquad
 |P^{(r)}(z)|\le C_{r,h}q^re^{hq}.
\]
A fixed-radius complex neighbourhood of $E_*(p_n/q_n)$ lies in $U_h^{\C}$ for all sufficiently large $n$.
\end{corollary}
\begin{proof}
Differentiate the product. Each derivative inserts $\left(\begin{smallmatrix}1&0\\0&0\end{smallmatrix}\right)$ at a distinct position. There are at most $q^r$ ordered choices and at most $r+1$ rational segments. Apply~\eqref{eq:complexrat} and bound the trace by twice the norm. The Chambers correction is independent of energy. Finally~\eqref{eq:continuity} gives convergence of the rational top edge to $\max\Sig$; compact containment in the open complex neighbourhood supplies the fixed radius.
\end{proof}

For real energies this is also exactly the derivative input (I2) of Paper~A, Proposition~3.3(2). Before using it on a whole rational band,~\eqref{eq:continuity} puts the whole phase union in its neighbourhood. It is never applied at an unanchored critical point far from the spectrum.

\label{end:complex}
